\documentclass[12pt]{article}
% AAU theme adapted from the preamble supplied by the project team.
\usepackage[a4paper,margin=1in]{geometry}
\usepackage{lmodern}
\usepackage[T1]{fontenc}
\usepackage{microtype}
\usepackage{graphicx}
\usepackage{float}
\usepackage{amsmath}
\usepackage[table]{xcolor}
\usepackage{titlesec}
\usepackage{fancyhdr}
\usepackage[most]{tcolorbox}
\PassOptionsToPackage{obeyspaces}{url}
\usepackage[hidelinks]{hyperref}
\usepackage{array}
\usepackage{tabularx}
\usepackage{enumitem}
\usepackage{booktabs}
% Vector UML artwork and multipage test tables.
\usepackage{longtable}
\usepackage{pdflscape}
\definecolor{aauRed}{RGB}{164,22,35}
\definecolor{aauGreen}{RGB}{0,122,88}
\definecolor{darkGray}{RGB}{60,60,60}
\definecolor{softGray}{RGB}{245,245,245}
\pagestyle{fancy}\fancyhf{}\fancyfoot[C]{\thepage}
\renewcommand{\headrulewidth}{0pt}\renewcommand{\footrulewidth}{0pt}
\titleformat{\section}{\normalfont\Large\bfseries\color{aauRed}}{\thesection}{0.6em}{}[{\color{aauRed!65}\titlerule[1.1pt]}]
\titlespacing*{\section}{0pt}{12pt}{8pt}
\titleformat{\subsection}{\normalfont\normalsize\bfseries\color{aauRed!85!black}}{\thesubsection}{0.6em}{}
\titlespacing*{\subsection}{0pt}{10pt}{4pt}
\newtcolorbox{objectivebox}{enhanced,colback=aauRed!5,colframe=aauRed!80!black,boxrule=0.3mm,arc=1mm,left=10pt,right=10pt,top=6pt,bottom=6pt,borderline west={2.5pt}{0pt}{aauRed}}
% Each short use/test case stays on one page, including its heading.
\newtcolorbox{contentbox}[1][]{enhanced,colback=softGray,colframe=black!35,boxrule=0.25mm,arc=2mm,left=9pt,right=9pt,top=7pt,bottom=7pt,before skip=9pt,after skip=9pt,title={\textbf{#1}},fonttitle=\bfseries,coltitle=darkGray,colbacktitle=black!8}
\newcommand{\blankline}{\rule{\linewidth}{0.35pt}}
\newcommand{\insertfigure}[3]{\begin{figure}[H]\centering\fbox{\parbox[c][#2][c]{0.88\textwidth}{\centering\textcolor{black!50}{\Large #1}}}\caption{#3}\end{figure}}
\newcommand{\screen}[3][0.98]{\begin{figure}[H]\centering\includegraphics[width=#1\linewidth]{../assets/#2}\caption{#3}\end{figure}}
% Give wide UML models a dedicated landscape page with readable labels.
\newenvironment{umlpage}{%
\newgeometry{margin=0.65in}\begin{landscape}\pagestyle{empty}%
\AddToHook{shipout/foreground}[arena/umlfooter]{%
\begin{tikzpicture}[remember picture,overlay]
\node[rotate=90,anchor=south,inner sep=0pt] at ([xshift=-23pt]current page.east)
{\normalfont\normalsize\thepage};
\end{tikzpicture}}%
}{\end{landscape}\RemoveFromHook{shipout/foreground}[arena/umlfooter]\pagestyle{fancy}\restoregeometry}
\newcommand{\umlfigure}[2]{\begin{figure}[H]\centering
\includegraphics[width=\linewidth,height=13cm,keepaspectratio]{../assets/uml/#1.pdf}
\caption{#2}\end{figure}}
\providecommand{\reportdate}{10 November 2026}
% Two intentionally sized figures fit together with a heading on an A4 page.
\newcommand{\screenpair}[5][0.30]{%
\begin{figure}[H]\centering
\includegraphics[width=0.96\linewidth,height=#1\textheight,keepaspectratio]{../assets/#2}\caption{#3}
\vspace{8pt}
\includegraphics[width=0.96\linewidth,height=#1\textheight,keepaspectratio]{../assets/#4}\caption{#5}
\end{figure}}
\setlength{\parindent}{0pt}\setlength{\parskip}{5pt}
\setlist[itemize]{leftmargin=1.5em,itemsep=2pt,topsep=3pt}
\setlist[enumerate]{leftmargin=1.7em,itemsep=3pt,topsep=4pt}
\setlength{\emergencystretch}{2em}
\widowpenalty=10000\clubpenalty=10000\displaywidowpenalty=10000
\brokenpenalty=10000
\raggedbottom
\renewcommand{\arraystretch}{1.16}
\setlength{\abovecaptionskip}{7pt}
\setlength{\belowcaptionskip}{4pt}
\setlength{\intextsep}{12pt plus 2pt minus 2pt}
\newcolumntype{Y}{>{\raggedright\arraybackslash}X}
\BeforeBeginEnvironment{tabularx}{\par\noindent}
\AfterEndEnvironment{tabularx}{\par}
\newcommand{\code}[1]{{\small\expandafter\nolinkurl\expandafter{\detokenize{#1}}}}
\newcommand{\command}[1]{\mbox{\small\ttfamily\detokenize{#1}}}
\newcommand{\projecttitle}[2]{%
\hypersetup{pageanchor=false}\begin{titlepage}\centering\vspace*{-0.45cm}
\noindent\begin{minipage}[c]{0.66\textwidth}\raggedright
{\Large\bfseries Al Ain University (AAU)\par}
\end{minipage}\hfill
\begin{minipage}[c]{0.28\textwidth}\raggedleft
% The supplied PNG is unchanged; trim removes only its transparent margins.
\includegraphics[width=3.1cm,trim=1731bp 871bp 1735bp 878bp,clip]{../../aau.png}
\end{minipage}\par
\vspace{0.45cm}
{\scshape\large Object Oriented Analysis and Design\par}
\vspace{0.10cm}
{\normalsize System Analysis \& Design\par}
\vspace{0.10cm}
{\small Course codes: 0103430 \;|\; 0102331\par}
\vspace{0.3cm}\noindent\rule{\linewidth}{0.5mm}\par\vspace{0.3cm}
{\huge\bfseries Arena\par}
\vspace{0.18cm}{\Large\color{aauRed}#1\par}
\vspace{0.15cm}{\normalsize\color{darkGray}#2\par}
\vspace{0.12cm}{\color{aauRed}\rule{\linewidth}{1.2pt}\par}
\vspace{0.35cm}{\normalsize\bfseries Prepared by\par}\vspace{0.15cm}
\begin{tabular}{@{}lr@{}}\textbf{Name}&\textbf{Student ID}\\\midrule
Khaled Alkhawaldeh&202410091\\Abdulrahman Alnaqbi&202340042\\Yousif Alhammadi&202310707\\Ahmed Alsamma&202320192\end{tabular}\par
\vspace{0.22cm}{\color{aauRed!55}\rule{0.6\textwidth}{1.2pt}\par}
\vspace{0.15cm}{\normalsize\textbf{Section:} 61D\par}
\vspace{0.15cm}{\normalsize\textbf{Submitted to:} Dr. Yazeed Ghadi\par}
\vspace{0.3cm}
\begin{tcolorbox}[colback=gray!5,colframe=black!70,title=\textbf{Project Summary},fonttitle=\bfseries,width=0.95\textwidth,boxrule=0.4mm,arc=3mm,left=7pt,right=7pt,top=4pt,bottom=4pt]
A Java desktop workspace for organizing campus competitions, coordinating teams, collecting submissions, evaluating work, and publishing results. Project, programming, and quiz scoring are separated from score-ranking, knockout, and league progression through object-oriented strategies. The application is developed for Al Ain University in the United Arab Emirates.
\end{tcolorbox}
\vfill{\normalsize Fall, 2026--2027\par}{\small Version 2.0 \;|\; \reportdate\par}
\end{titlepage}\pagenumbering{roman}\hypersetup{pageanchor=true}\tableofcontents\clearpage\pagenumbering{arabic}}

\begin{document}
\projecttitle{Implementation Report}{Delivered features, validation, and user guide}
\section{Introduction}
\subsection{Project Overview}
Arena is a Java desktop application for university competitions. It addresses the instructor's request for a project that goes beyond basic registration: project showcases use multi-judge weighted rubrics, programming contests use solved-count and penalty ranking, quiz competitions record bounded official points, and separate knockout/round-robin strategies generate matches and determine champions.
\subsection{Project Objectives}
The delivered objectives are a usable desktop interface, persistent local data, four authenticated roles, validated team workflows, format-aware evaluation, protected result publication, an executable test suite, and a coherent set of analysis, design, testing, and interface documents.

\section{Implementation}
\subsection{Development Environment and Technologies}
The source targets Java 17 and was compiled on Windows using a locally downloaded JDK 21 with \command{--release 17}. The downloaded JDK archive's SHA-256 matched its release metadata. The app uses Swing/AWT, collections and records, Java time, Java cryptography, and NIO filesystem APIs. No external runtime library or database server is necessary.

The project includes \code{build-app.ps1}, \code{start-app.ps1}, a Windows double-click launcher, a POSIX build script, and an optional Maven descriptor. The executed validation command was \command{./build-app.ps1 -Test -UiTest}. Maven also executes the same suites through JUnit and Surefire. Use mvn test for service checks and mvn test -Pui-tests for both suites on a graphical desktop.
\subsection{Implemented Features}
\begin{tabularx}{\linewidth}{lY}\toprule Area & Delivered behavior\\\midrule
Accounts & Sign in/out, role assignment through administrator account creation, own-password change.\\
Competitions & Create/edit draft rules, assign judges, browse/search/filter, enforce staged workflow.\\
Teams & Create, self-join, leave as noncaptain, show roster and captain, lock registered membership.\\
Entries & Captain registration/withdrawal with duplicate, overlap, size, capacity, and phase checks.\\
Submissions & Title, description, valid HTTP(S) link, deadline enforcement, replacement before judging, scoped read access.\\
Project judging & Three criteria weighted 40/35/25; independent assigned judges; average and shared ranks.\\
Programming & Official solved count and penalty; exclusive score owner for each entry.\\
Quiz & One official bounded points result per team; ties share rank; quiz match formats use both teams' points.\\
Knockout & Seeded first round with byes for 2--16 teams, immutable match results, automatic rounds, champion certification.\\
Round robin & Every pair once; 3/1/0 league points; complete standings, deterministic tie-breaks, certification.\\
Results & Manager previews, publication gates, leaderboard, safe CSV export, archive.\\
Persistence & Local snapshot, process lock, filtered reads, save rollback, audit history.\\\bottomrule\end{tabularx}
\subsection{Implementation Challenges and Solutions}
\textbf{Fairness across formats.} A generic score could not represent both judging and elimination. The solution separates CompetitionType from TournamentFormat. CompetitionRules handles scoring; TournamentRules handles fixture generation, progression, and match standings. The service enforces compatibility and permissions for both.

\textbf{Incomplete publication.} A partial leaderboard could misrepresent unreviewed teams. Publication checks the complete expected evaluation set, while provisional views remain available only to competition managers.

\textbf{Registration consistency.} A student could otherwise enter twice by joining another team after registration. Roster locking and overlapping-member validation close that path.

\textbf{Persistence failures.} Mutating memory before a failed save could make the interface show data that disappears after restart. Each mutation now has a rollback snapshot. A deliberately blocked temporary-file write was tested to verify restoration.

\textbf{Local build portability.} The starting environment did not expose a usable Java compiler. A portable JDK was placed in the ignored tools directory. Scripts also support an installed JDK through JAVA\_HOME or PATH. The JAR tool's temporary directory was directed into the build folder after the sandbox rejected its default cross-directory replacement.

\textbf{Screen quality.} Actual Swing renders were inspected. An introductory paragraph was clipped because its preferred height was too small; the layout was corrected before regenerating the documentation images.

\section{Design Changes}
\subsection{Deviations from the Original Design}
The repository initially contained task outlines, not an implemented baseline architecture. The following are recorded implementation decisions rather than claims of a previously approved design being replaced:
\begin{itemize}
\item Use a local Swing desktop app to meet the Java requirement without a web deployment or external database installation.
\item Keep roles in one immutable User record rather than inventing subclasses with identical fields.
\item Separate UI, services, repository, rules, utilities, and account summaries into packages. Preserve serialized Model record names for existing-data compatibility. Extract dialogs and competition views from ArenaApp.
\item Match the clickable prototype with burgundy actions, a white sidebar, serif headings, a feature banner, competition cards, and an assigned Judging studio.
\item Present judging through a submission reader and scoring dialog instead of the wireframe's permanent two-column pane.
\item Support deadlines and bracket rounds while deferring a full venue/time-slot scheduling calendar.
\end{itemize}
\subsection{Reasons for Changes}
These choices keep the application easy to run and demonstrate while preserving substantive OOAD behavior. They reduce setup costs and unnecessary inheritance, but impose explicit limitations: one process, whole-state persistence, fixed scoring policies, and manual sequential sign-in for different actors.

\section{Testing Summary}
\subsection{Testing Approach}
The dependency-free Java suite exercises application services with a fixed clock and isolated temporary data. It combines formula checks, negative authorization and input cases, end-to-end format workflows, persistence reopen checks, lock contention, corrupt-data handling, and injected save failure. ArenaTournamentTest adds 28 regression checks, including exhaustive valid tournament sizes. A Swing runner exercises 22 real control and dialog scenarios on the event-dispatch thread. Actual Swing screens are also rendered for visual inspection.
\subsection{Results and Quality Assessment}
\begin{objectivebox}
Executed validation: \textbf{98 service/regression checks and 22 Swing UI checks passed; 0 failed} on 10 November 2026. The machine-readable results and execution summary are preserved under \code{reports/test-evidence}. These results describe the executed suite; they do not imply complete branch coverage, production security certification, or a user study.
\end{objectivebox}
The detailed Testing Report maps each test to its requirement and documents expected and actual outcomes. Thirteen application screens and dialogs were rendered. Seventeen browser prototype checks also passed. The desktop app and browser prototype share the same palette, type hierarchy, card layout, and judging terminology. Windows is the verified platform; high-DPI, screen-reader, multi-platform, and large-dataset validation remain outstanding.

\section{Project Status}
\subsection{Completed Features}
The source, runnable JAR build, seeded demonstration dataset, automated suite, analysis and design models, interface wireframes, screenshot assets, LaTeX reports, user instructions, and demo walkthrough are included. The original task outlines are preserved for comparison.
\subsection{Incomplete Features and Known Limitations}
\begin{itemize}
\item This is a single-computer coursework application, not a hosted campus service. No concurrent client sessions, SSO, email, or remote database is implemented.
\item Programming scores are entered by a judge; there is no automatic code grader or sandbox.
\item Knockout byes and round-robin leagues are implemented for 2--16 teams. Double elimination, fixture time slots, cancellation, and post-result correction remain outside scope.
\item Every project/programming ranking entry must submit before judging begins; non-submission can stall those events. Quiz and match formats require no project submission. Forfeit/disqualification and cancellation flows are future work.
\item Submission links are stored, not verified remotely. Files are not uploaded.
\item No roster invitations, captain transfer, team deletion, account deletion, or credential-reset workflow is provided. Once registered, a roster is intentionally locked.
\item Version-one records migrate to the Arena type/format model in memory and persist as schema two on the next successful save. Both completed and active legacy fixtures are tested. Manual backup while closed is required; data is not encrypted.
\end{itemize}

\clearpage\section{Team Collaboration}
\subsection{Team Responsibilities}
The names and IDs below were supplied by the team. The responsibility allocation is a \textbf{proposed presentation and review plan}, not a record of work independently performed by each person.

\begin{tabularx}{\linewidth}{lY}\toprule Member & Proposed review/presentation responsibility\\\midrule
Khaled Alkhawaldeh & Integration, requirements traceability, and overall demonstration.\\
Abdulrahman Alnaqbi & Team/registration rules and user-guide walkthrough.\\
Yousif Alhammadi & Format strategies, scoring policies, and bracket explanation.\\
Ahmed Alsamma & Test execution, persistence behavior, and document review.\\\bottomrule\end{tabularx}

The implementation and documents were prepared with AI assistance at the team's request. Actual individual contributions remain unverified and require group confirmation before submission. The group should review, understand, and adapt the work and follow its course's disclosure requirements.
\subsection{GitHub Development Workflow}
Changes are made in the local repository. Committing and pushing are left to the team as requested; no repository history, commits, or pushes are fabricated. A suggested workflow is to review the diff, run the tests, review the LaTeX PDFs, commit source and documents, and push through the team's usual GitHub process. Runtime downloads, private local databases, and build outputs are ignored. The initial Git history had four commits and tracked only the README, ignore file, and five task outlines; application and report files were outside that snapshot. The final submission contains the complete current source and documents. No historical authorship is inferred from this snapshot.

\section{Conclusion}
The delivered system makes competition variation visible in both the domain design and executable behavior. It integrates registration with submissions, judging, bracket progression, and guarded publication, supported by reproducible checks and consistent documentation. The next major architectural step would be a database-backed service with separate authenticated clients, followed by richer tournament scheduling and recovery workflows.
\clearpage\section{References}
\begingroup\raggedright\begin{enumerate}
\item Object Management Group. \emph{Unified Modeling Language, version 2.5.1}. \url{https://www.omg.org/spec/UML/2.5.1}. Accessed 29 October 2026. Used as the notation reference for the UML models.
\item Oracle. \emph{Java SE 17: javax.swing package}. \url{https://docs.oracle.com/en/java/javase/17/docs/api/java.desktop/javax/swing/package-summary.html}. Accessed 29 October 2026. Used for desktop components and event-dispatch-thread conventions.
\item Oracle. \emph{Java SE 17: ObjectInputFilter}. \url{https://docs.oracle.com/en/java/javase/17/docs/api/java.base/java/io/ObjectInputFilter.html}. Accessed 29 October 2026. Used for bounded, filtered repository deserialization.
\item Project materials. Task outlines in \code{requirements/system-analysis-requirements.txt}, \code{requirements/system-design-requirements.txt}, \code{requirements/implementation-report-requirements.txt}, \code{requirements/testing-requirements.txt}, and \code{requirements/sample-data-requirements.txt}; instructor feedback supplied by the team. These define the coursework scope, not an independent evaluation of the implementation.
\end{enumerate}\endgroup


\clearpage\appendix
\section{User Guide}
\subsection{System Requirements}
Windows with PowerShell, a graphical desktop, and JDK 17 or newer. The local tools directory provides the JDK used in this workspace but is not committed. On another machine, install a JDK and configure JAVA\_HOME or PATH. The minimum application window is 1100 by 720; the default is 1440 by 1000.
\subsection{Installation Instructions}
Clone or copy the project folder. No runtime dependencies need to be downloaded by the build. If using an IDE, import the Maven project or open the Java source directory. The supplied PowerShell build resolves JAVA\_HOME first, then the local portable JDK, then PATH.
\subsection{Running the Application}
On Windows, double-click \code{start-app.cmd}, or use:
\begin{verbatim}
.\start-app.ps1
.\build-app.ps1 -Test
java -jar build/arena.jar
# Maven alternative (downloads build/test dependencies):
mvn clean package -Pui-tests
java -jar target/arena.jar
\end{verbatim}
A fresh launch creates \code{data/arena.bin}. If \code{data/campus.bin} already exists, the launcher and Java entry point reuse it, preserving the existing workspace. To use a separate demo database, run:
\begin{verbatim}
.\start-app.ps1 -DataFile data/demo.bin
\end{verbatim}
On POSIX systems, use \command{sh build-app.sh test} and then the Java JAR command; the script uses LF endings and supports the Windows Git shell as well as a POSIX shell. Native Linux/macOS execution is unverified.
\subsection{User Roles and Permissions}
All seeded accounts initially use \code{Campus@2026}. Organizer: \code{organizer1@campus.test}; administrator: \code{admin@campus.test}; judges: \code{judge1@campus.test} through \code{judge3@campus.test}; students: \code{student1@campus.test} through \code{student10@campus.test}. A second organizer is also available. These are fictional local demo identities, not deliverable email addresses.

Student1 captains Team Alpha; student3 captains Code Warriors; student5 captains Innovators; student7 captains Circuit Breakers; student9 captains Nova Labs. Each has the following even-numbered student as a member.
\subsection{Step-by-Step Usage Instructions}
\textbf{Project showcase demonstration:}
\begin{enumerate}
\item Sign in as organizer1 and open Hackathon 2026. Three teams are already registered; judges 1 and 2 are assigned. Click Start competition.
\item Sign out. Sign in as student1 and submit Team Alpha's title, description, and project link. Repeat as student3 and student5 for their teams.
\item As organizer1, choose Begin judging. Missing submissions produce a clear rejection rather than silently exclude a team.
\item As judge1, read each submission and record innovation, execution, presentation, and feedback. Repeat all three reviews as judge2.
\item As organizer1, preview the results and choose Publish results. As a student, open Results and export the official CSV if needed.
\end{enumerate}
\textbf{Programming demonstration:} the seeded Programming Contest is already published. Code Warriors leads with nine solved problems. Team Alpha and Innovators both solved eight; Team Alpha ranks higher because its penalty is lower.

\textbf{Knockout demonstration:} sign in as judge3, open Robotics Challenge, and select Bracket. Record each semifinal winner. The final appears when both semifinals are complete. Record its winner; sign in as organizer1 to certify the bracket and publish results.

\textbf{Quiz ranking demonstration:} as judge2, open Academic Quiz and record points from 0--20 for each of the three registered teams. As organizer1, publish after all official results exist. Equal points share ranks. No submission is required.

\textbf{League demonstration:} as judge3, open Campus League and select Fixtures. Select each of the three fixtures, record both scores, and save. Draws award one point each. As organizer1, inspect Results preview, certify the completed fixtures, and publish. The CSV includes all league statistics.

\textbf{Creating an event:} as organizer1, choose Create competition, select the competition type, select one of its compatible tournament formats, configure quiz maximum points if applicable, and Ctrl-click the intended judges. The form explains the rubric or match rules and adjusts the capacity limit. Use an ISO local deadline such as \code{2026-12-01T17:00}. Open registration. Student captains register through their own accounts. Rules cannot be edited after opening.
\clearpage\subsection{Screenshots and Examples}
\screenpair{screenshots/02-overview.png}{Organizer overview after sign-in.}{screenshots/06-results.png}{Official sample programming results.}
The Interface Mock-ups document contains thirteen application screen/dialog captures, seven browser-prototype captures, and four editable skeletons.
\subsection{Troubleshooting}
\begin{tabularx}{\linewidth}{lY}\toprule Symptom & Resolution\\\midrule
Compiler not found & Install JDK 17+ and set JAVA\_HOME; a JRE alone does not compile sources.\\
Script execution blocked & Use the provided Windows launcher or invoke PowerShell with the script-specific execution-policy flag; do not change the machine-wide policy.\\
Data already open & Close the other Arena process using the same data file, or select another data path.\\
Cannot advance & Read the guard message: verify entry count, all submissions, reviews, or bracket completion.\\
Cannot register team & Check captain identity, phase, deadline, capacity, roster size, and overlapping members.\\
Cannot edit roster & Registered rosters are locked. Withdraw while registration is open, if appropriate, or use another unregistered team.\\
Corrupt data error & Preserve the file; restore a known backup while the app is closed. Do not overwrite it to hide the error.\\
Expired demo deadlines & Existing demo data does not refresh itself. Create a new event or use a fresh, separately named data file.\\\bottomrule\end{tabularx}
\end{document}
